% 共存試験: 行列・場合分けのセル内和文（CJKアクティブ文字と\halignの相互作用）。
% root cause（修正済み）: backend の post-glyph peek \jlreqpdf@peek（=
% \CJK@ignorespaces）が「\ifmmode\else\futurelet..\fi」形だったため、CJKutf8の
% \CJK@XXX 末尾条件群内で展開されると \else スキップを越えて窥き見本体が古い
% \jlreqpdf@next で走り、\CJK@XXX の \fi 群を突いて Extra \fi で落ちた（再現:
% 直前に和文がある状態の $あ$、行列セル内和文）。\@gobble イディオム化＋数式内
% グリフの自然幅化で解消。以下がすべて組めることを検証する。
\documentclass[pdflatex]{jlreq}
\usepackage{amsmath}
\begin{document}
行内: \(\begin{pmatrix} あ & b \\ c & d \end{pmatrix}\) と、
行列内和文を含む式の前後「前 \(\begin{pmatrix} あ \\ い \end{pmatrix}\) 後」。
別行立て:
\[
\begin{pmatrix}
\mathmc{漢} & あ & x \\
y & \text{和} & 1
\end{pmatrix}
\quad
\begin{bmatrix}
\alpha & （あ） \\
い & \beta
\end{bmatrix}
\quad
\begin{Bmatrix}
。 & 、 \\
! & ?
\end{Bmatrix}
\]
場合分け: \(f(x)=\begin{cases} あ & x>0 \\ い & x\le 0 \end{cases}\)。
直前が和文で数式: 解答は \(あ=1\) である。次の文、\(い=2\) も確認。
\end{document}
